\documentclass[../main]{subfiles}
\begin{document}
\chapter{Corriente Alterna}
\img{images/11/alterna}{Corriente Alterna.}

\info{Corriente Alterna}{
Se denomina corriente alterna (AC, por sus siglas en inglés de Alternating current) a la corriente eléctrica en la que la magnitud y el sentido varían cíclicamente.

La forma de oscilación de la corriente alterna más comúnmente utilizada es la oscilación senoidal con la que se consigue una transmisión más eficiente de la energía, a tal punto que al hablar de corriente alterna se sobrentiende que se refiere a la corriente alterna senoidal.
}

\info{Velocidad angular}{
Velocidad angular: Es la velocidad angular en radianes que un fasor de corriente alterna.

\begin{equation*}
    \omega = 2 \times \pi \times f 
\end{equation*}

Forma de la onda
\begin{equation*}
    v = V_{MAX} \times  sen(\omega \times t+\varphi)
\end{equation*}

\begin{equation*}
    i = I_{MAX} \times  sen(\omega \times t+\varphi)
\end{equation*}

Siendo:

$v$ y $V_{MAX}$: f.e.m. de la onda sinusoidal instantánea, y valor máximo de la onda sinusoidal en volts $[V]$.

$i$ y $I_{MAX}$: corriente de la onda sinusoidal instantánea, y valor máximo de la onda sinusoidal en Amperes $[A]$.

$\omega$: velocidad angular en radianes por segundo $[rad/s]$

$\varphi$: ángulo de atraso o adelanto de onda en radianes [rad]

$t$: tiempo en segundos $[s]$.
}

\info{Frecuencia y período}{
La frecuencia se mide en Herz $[Hz]$ Cuya unidad es equivalente a $[1/s]$ y nos da una idea de los ciclos que ocurren en 1 segundo.
\begin{equation*}
    f=\frac{1}{T} \Rightarrow T=\frac{1}{f} \Rightarrow f\times T = 1
\end{equation*}




El período $T$, es el tiempo que demora en ocurrir el ciclo y se mida en segundos [s].
\begin{equation*}
1000 ms (milisegundos) = 1s
\end{equation*}

}

\info{Valor Medio y Eficaz de una tensión o corriente alterna.}{
\begin{equation*}
    V_m= \frac{2\times V_{MAX}}{\pi} = 0.636 \times V_{MAX}
\end{equation*}

\begin{equation*}
    I_m= \frac{2\times I_{MAX}}{\pi} =  0.636 \times I_{MAX}
\end{equation*}

\begin{equation*}
    V= \frac{V_{MAX}}{\sqrt{2}} = 0,707 \times V_{MAX}
\end{equation*}

\begin{equation*}
    I= \frac{I_{MAX}}{\sqrt{2}} = 0,707 \times I_{MAX}
\end{equation*}

}
\actividad{Resolver los siguientes problemas de corriente alterna:}
\begin{enumerate}
    \item Hallar el valor de la tensión instantánea $v$ de un alternador con frecuencia $f=50Hz$ y $V_{MAX}=250V$ con $\varphi=0°$ y $\omega=1 \frac{rad}{s}$ después de $t=1,4s$ segundos de tiempo.
    \item ¿Cuál el valor de la tensión instantánea $v$ de un alternador con frecuencia $f=60Hz$ y $V_{MAX}=350V$ con $\varphi=\pi$ y $\omega=1,5 \frac{rad}{s}$ después de $t=10s$ segundos de tiempo?
    
    \item Calcular las frecuencias de corrientes alternas que tienen períodos correspondientes: $T_1=20ms$, $T_2=40ms$ y $T_3=1ms$
    \item Calcular el período de cada una de las f.e.m alternas cuyas frecuencias son: $f_1=200Hz$, $f_2=1000Hz$ y $f_3=50Hz$.
    \item ¿Cuál es el valor máximo $V_{MAX}$ y medio $V_m$ de una tensión de valor eficaz $V=220V$?
    \item Calcular el valor máximo $I_{MAX}$ y medio $I_m$ de una corriente de valor eficaz $I=10A$.
    \item ¿Cuál es el valor máximo $V_{MAX}$ y medio $V_m$ de una tensión de valor eficaz $V=380V$?

    \item Hallar el valor de la tensión instantánea $v$ de un alternador con frecuencia $f=50\,\text{Hz}$ y $V_{\text{MAX}}=150\,\text{V}$ con $\varphi=\frac{\pi}{4}$ y $\omega=2\,\text{rad/s}$ después de $t=2,5\,\text{s}$ segundos de tiempo.
    
    \item ¿Cuál es el valor de la tensión instantánea $v$ de un alternador con frecuencia $f=75\,\text{Hz}$ y $V_{\text{MAX}}=300\,\text{V}$ con $\varphi=\frac{\pi}{3}$ y $\omega=2,5\,\text{rad/s}$ después de $t=5\,\text{s}$ segundos de tiempo?
    
    \item Calcular las frecuencias de corrientes alternas que tienen períodos correspondientes: $T_1=10\,\text{ms}$, $T_2=25\,\text{ms}$ y $T_3=2\,\text{ms}$.
    
    \item Calcular el período de cada una de las f.e.m alternas cuyas frecuencias son: $f_1=60\,\text{Hz}$, $f_2=500\,\text{Hz}$ y $f_3=200\,\text{Hz}$.
    
    \item Hallar el valor de la corriente instantánea $i$ para una corriente alterna con frecuencia $f=50\,\text{Hz}$, $I_{\text{MAX}}=20\,\text{A}$, y $\varphi=0$ después de $t=0,1\,\text{s}$.
    
    \item Calcular el valor eficaz $I_{\text{eficaz}}$ de una corriente cuya forma de onda es $I(t) = 15 \sin(\omega t + \frac{\pi}{2})$, donde $I_{\text{MAX}} = 15\,\text{A}$ y $f = 50\,\text{Hz}$.
    
    \item Calcular el período y la frecuencia de una onda alterna con $\omega = 314\,\text{rad/s}$.
    
    \item Hallar el valor instantáneo de la tensión $v$ para una onda alterna con $V_{\text{MAX}} = 400\,\text{V}$, $f = 60\,\text{Hz}$, $\varphi = \frac{\pi}{6}$, después de $t = 3,2\,\text{s}$.
    
    \item ¿Cuál es el valor medio $I_m$ y el valor eficaz $I_{\text{eficaz}}$ de una corriente alterna con valor máximo $I_{\text{MAX}} = 25\,\text{A}$?

\img{images/11/senoidal}{Partes de la onda seno.}

    
    \item Calcular la frecuencia y el período de una corriente alterna si la velocidad angular $\omega = 628\,\text{rad/s}$.

    \item Hallar la potencia instantánea $P(t)$ de un circuito donde la tensión es $v(t) = 220\sqrt{2} \sin(100\pi t)$ y la corriente es $i(t) = 10\sqrt{2} \sin(100\pi t + \frac{\pi}{4})$.
    
    \item Calcular la energía total disipada en una resistencia de $50\,\Omega$ durante un ciclo completo de una corriente alterna $i(t) = 5\sin(314t)$.
    
    \item Un generador de corriente alterna produce una tensión de $v(t) = 300\cos(377t)$. Hallar la energía almacenada en un condensador de $10\,\mu F$ conectado al generador después de $t=0,005\,\text{s}$.
    
    \item ¿Cuál es el valor eficaz de una corriente alterna cuya forma de onda es $i(t) = 5\sin(100t) + 3\sin(300t)$? 

   
\end{enumerate}


\end{document}